\subsection{Координаты: где точка}

Точки задаются в \textbf{локальной системе координат} в \textbf{метрах}:
\begin{itemize}
    \item ось \textbf{X} направлена вперёд (туда, куда смотрит нос дрона на старте), ось \textbf{Y} --- влево, ось \textbf{Z} --- вверх;
    \item начало координат $(0, 0, 0)$ --- стартовая площадка. Это фиксированная точка сцены, а не место последнего взлёта: \verb|ap.goToLocalPoint(0, 0, 1)| всегда возвращает дрон в точку над стартом;
    \item на настоящем Пионере координаты даёт система позиционирования (ультразвуковая LPS или оптическая); в симуляторе --- сцена.
\end{itemize}

\begin{center}
\simshot[0.85\linewidth]{06-coords.png}
\captionof{figure}{Камера \ui{Земля} --- вид сверху, оси как на компасе в правом верхнем углу. \key{Ctrl}+щелчок по земле или объекту показывает координаты точки, которые можно подставить в \texttt{ap.goToLocalPoint}.}
\end{center}

\begin{hintbox}[Z --- это высота цели, а не поверхности]
\key{Ctrl}+щелчок показывает высоту поверхности, по которой щёлкнули. Лететь нужно выше: над крышей высотой 13~м задавайте $z = 14$. Над ровной площадкой удобная учебная высота --- от 0.7 до 1.5~м.
\end{hintbox}

\subsubsection*{Курс}

Курс (\texttt{yaw}) --- поворот носа вокруг вертикальной оси. Автопилот принимает его в \textbf{радианах}: $\pi/2 \approx 1.57$ --- четверть оборота влево. Переводить градусы удобно функцией \verb|math.rad(deg)|. Команда \verb|ap.updateYaw(rad)| меняет курс, не меняя точку; её часто вызывают перед \verb|goToLocalPoint|, чтобы дрон летел носом вперёд.

\begin{taskbox}[Разведка площадки]
Включите камеру \ui{Земля} и \key{Ctrl}+щёлкните по четырём точкам вокруг старта. Запишите их координаты --- они понадобятся для маршрута в следующем разделе.
\end{taskbox}
